\cleardoublepage
\phantomsection
\addcontentsline{toc}{section}{Заключение}
\section*{Заключение}

В работе получены количественные оценки достаточного размера обучающей выборки параметрических моделей на основе анализа стабилизации поверхности эмпирической функции потерь в окрестности ее минимума. Сформулированные во введении задачи решены следующим образом:

\begin{enumerate}
    \item Проведен анализ существующих подходов к измерению стабилизации ландшафта (глава~\ref{sec1}); показано, что одноточечная оценка и изотропное среднеквадратичное усреднение в полном пространстве параметров являются частными случаями более общей конструкции, причем доминирующая константа в верхней оценке последнего критерия растет с размерностью пространства параметров как $N^2$, что делает его количественное применение в нейронных сетях современных размеров невозможным.
    \item Введен унифицированный критерий $p$-сходимости $\Delta_p$ с функцией предпочтения точек в пространстве параметров (глава~\ref{sec2}); получено разложение разности эмпирических функций потерь в окрестности минимума на три вклада~\eqref{eq:losses-difference-second-order-minimum}.
    \item Доказаны оценки скорости убывания критериев: $\mathcal O(k^{-1})$ для $\Delta_1$, $\mathcal O(k^{-2})$ для $\Delta_2$ и $\Delta_2^{(D)}$ (теоремы~\ref{thm:abs-onepoint-general},~\ref{thm:squared-criterion-general},~\ref{thm:squared-subspace-general}); установлены замкнутое спектральное представление и экстремальность подпространства главных направлений (теорема~\ref{thm:squared-subspace-hessian-eigs}, утверждение~\ref{prop:extremal-topD}); построены три масштабируемых алгоритма оценивания критерия (раздел~\ref{subsec:algorithms}). На основе оценок получено количественное выражение~\eqref{eq:sufficient-sample-size} для достаточного размера выборки.
    \item Полученные оценки эмпирически проверены (глава~\ref{sec4}): на MLP/MNIST показатели степенного убывания согласуются с теоретическим $\alpha=2$ (таблица~\ref{tab:landscape-exponents}); на трансформере \texttt{nanochat d6} подпространственный критерий воспроизводит полнопространственный сигнал при $D/N < 10^{-6}$ в локальном режиме $\sigma \leqslant 10^{-3}$, а оценка замкнутой формы в этом режиме приблизительно в $1{,}8\cdot 10^4$ раз быстрее прямой Монте-Карло без потери точности.
\end{enumerate}

Подробные формулировки научной новизны, положений на защиту, теоретической и практической значимости приведены во введении и не повторяются здесь.

\textbf{Ограничения работы.} Все полученные оценки опираются на локальную квадратичную аппроксимацию функции потерь, точность которой эмпирически подтверждена для $\sigma \leqslant 10^{-3}$. Утверждение~\ref{prop:extremal-topD} об экстремальности подпространства ведущих собственных векторов доказано в идеализированном предположении одновременной диагонализуемости матриц $\mathbf H^{(k)}(\mathbf w_k^*)$ и $\mathbf H^{(k+1)}(\mathbf w_k^*)$. Утверждение~\ref{prop:fcnn-G-bound} об архитектурно-зависимой оценке $M_{\mathbf H}$ сформулировано только для полносвязных сетей без смещений.

\textbf{Направления дальнейших исследований.} Перенос результатов на режимы с медленным дрейфом главных собственных направлений с использованием адаптивных пробных подпространств; связь количественной оценки~\eqref{eq:sufficient-sample-size} с эмпирическими законами масштабирования параметрических моделей; расширение архитектурно-зависимой оценки $M_{\mathbf H}$ на сверточные и трансформерные архитектуры.
